\documentclass[12pt,letterpaper]{hmcpset}
\usepackage{amsmath}
\usepackage{amssymb}
\usepackage{amsfonts}
\usepackage{graphicx} % Required for inserting images
\usepackage{hyperref}
\usepackage{enumitem}
\usepackage[margin=1in]{geometry}
\usepackage{graphicx}
\usepackage{float}
\DeclareMathOperator{\vspan}{span}
\newcommand{\bb}[1]{\mathbb{#1}}
\newcommand{\bvec}[1]{\mathbf{#1}}
\newcommand{\pardx}[2]{\frac{\partial #1}{\partial #2}}
\newcommand{\ddx}[2]{\frac{d #1}{d #2}}
\newcommand{\norm}[1]{\Vert #1 \Vert}
\newcommand{\abs}[1]{\vert #1 \vert}
\newcommand{\set}[1]{\{ #1 \}}
\setlist[enumerate, 1]{label = (\alph*)}

\name{Jayden Thadani}
\class{Principles of Mathematical Analysis self-study}
\assignment{Chapter 2 exericses---Basic Topology}

\begin{document}

\begin{problem}[16.]
	Regard $\bb{Q}$, the set of all rational numbers, as a metric space, with \( d(p, q) = \abs{p - q} \). Let $E$ be the set of all \( p \in \bb{Q} \) such that \( 2 < p^2 < 3 \). Show that $E$ is closed and bounded in $\bb{Q}$ but that $E$ is not compact. Is $E$ open in $\bb{Q}$?
\end{problem}

\begin{solution}
	First we will show that $E$ is bounded.\\\\
	Consider the point $q  = 1.5$ and the neighbourhood around it of radius $0.5$. We can show that $E$ is contained in that neighbourhood, and is thus bounded. We prove this by proving that any point not contained in the neighbourhood is not contained in the set as below.\\\\
	If \( d(p, q) > 0.5 \) then we have
	\[
		\abs{p - 1,5} > 0.5
	\]
	which gives us
	\( p - 1.5 > 0.5 \) or \( p - 1.5 < -0.5 \).\\\\
	Then we must have \( p > 2 \) or \( p < 1\) and thus \( p^2 < 1 \) or \( p^2 > 4 \). This gives us \( p \notin E \). Thus, we've shown that  $E$ is bounded.
	Now we will show that $E$ is closed by proving that any point that is not in $E$ cannot be a limit point of $E$.\\\\
	If \( p \notin E\) then we must have \( p^2 < 2 \) or \( p^2 > 3 \). In each of these cases e can construct a neighbourhood around $p$ that does not include any points of $E$.\\\\
	By the construction in chapter 1 to prove that the rationals do not have the least upper bound property, we know that there exists a positive rational $r$ such that for any $p$ with \( p^2 < 2 \) we have \( (p + r)^2 < 2 \). Specifically,
	\[
		r = \min{(\frac{2 - p^2}{2p + 1}, 1)}
	\]
	gives us this property. Moreover, since $r$ is positive and the square is an increasing function for positive reals, this gives us \( q^2 < 2 \) for any \( q \in (p - r, p + r) \). Thus $N_r(p)$ contains no points of $E$ making $p$ not a limit point of $E$ for any $p$ with \( p^2 < 2 \).\\\\
	We have a similar proof to show that $p$ with \( p^2 > 3 \) also cannot be a limit point of $E$. Particularly, we use a similar construction for $r$ such that the $N_r(p)$ contains no points of $E$ whenever \( p^2 > 3 \).\\\\
	Consider 
	\[
		r = \min{(\frac{p^2 - 3}{2p + 1}, 1)}
	\]
	for some $p$ with \( p^2 > 3 \).
	Then
	\[
		(p - r)^2 = p^2 - r(2p - r) > p^2 - r(2p + 1)
	\]
	since \( -r < 1 \). Substituting $r$ gives us
	\[
		(p - r)^2 > p^2 - (\frac{p^2 - 3}{2p + 1})(2p + 1) = p^2 - (p^2 - 3) = 3.
	\]
	That is
	\[
		(p - r)^2 > 3.
	\] \\
	Again, since $r$ is positive, and the square function is increasing, this gives us this gives us \( q^2 > 3 \) for any \( q \in (p - r, p + r) \). Thus $N_r(p)$ contains no points of $E$ making $p$ not a limit point of $E$ for any $p$ with \( p^2 > 3 \).\\\\
	Since we've exhausted all the cases for rationals not in $E$ (there is no rational $p$ with \( p^2 = 2 \) or \( p^2 = 3 \)) and shown all of them not to be limit points, we've shown that all the limit points of $E$ must be in $E$, and thus, $E$ must be closed.\\\\
	Finally, now we will show that $E$ is compact. Consider the cover of $E$ given by $\set{G_p}$ where \( G_p = (p - r_p, p + r_p) \) such that \( [p - r_p, p + r_p] \subseteq E \) for all \( p \in E \). Note that such an $r_p$ exists for each \( p \in E \) by a similar construction to that which we used earlier.\\\\
For any finite subset of $\set{G_p}$, say $S$, there must be some maximum $p + r_p$ chosen among all the $G_p$ in $S$. Then that maximum \( p + r \notin \bigcup_{G_p \in S} G_p \) but \( p + r \in E \) making $S$ not a finite subcover, and making $E$ not compact.
\end{solution}

\begin{problem}[17.]
	Let $E$ be the set of all \( x \in [0, 1] \) whose decimal expansion contains only the digits $4$and $7$. Is $E$ countable? Is $E$ dense in $[0, 1]$? Is $E$ compact? Is $E$ perfect?
\end{problem}

\end{document}

